Assume \(n=p\). Full row rank makes \(C\) square and invertible, so any basis extension with observed rows \(C\) is simply \(T_0=C\), with \(W_0=C^{-1}\). The Lacerda--Spirtes--Ramsey--Hoyer construction cited in the two supporting leaves permutes rows of this square unmixing matrix so that the diagonal is nonzero and then rescales rows to make that diagonal one; their Theorem 1 concerns distribution entailment, not intervention solvability.

Consider first any such candidate whose structural matrix \(Q\) also satisfies \(\det(Q_{-x,-x})\ne0\). Its row normalization has monomial source assignment, its mixing matrix is \(C\), and it satisfies all conditions of \(\mathfrak F_x(C,\varepsilon)\). The sharp-effect-frontier theorem therefore places its effect in \(\mathcal R_{xy}(C)\).

For the reverse inclusion, take \(j\in J_x(C)\). In the square case \(\mathsf L=\varnothing\). Criterion (5) from the shear-and-matching proof specializes to
\[
 \operatorname{rank}(C_{-x,-j})<p-1
 \quad\Longleftrightarrow\quad (C^{-1})_{jx}=0.
\]
Thus admissibility implies \((C^{-1})_{jx}\ne0\), and no latent-row shear is needed. The nonsingular complementary minor \((C^{-1})_{-j,-x}\) supplies a perfect matching; augmenting it by \(j\mapsto x\), permuting the rows, and normalizing their diagonal is exactly a square admissible row-permutation candidate of the cited construction. The calculation in the sufficiency lemma proves that its \(Q_{-x,-x}\) is nonsingular and that its effect is \(C_{yj}/C_{xj}\). Every element of \(\mathcal R_{xy}(C)\) is therefore attained by a postintervention-solvable square candidate, proving the claimed equality. Candidates with singular \(Q_{-x,-x}\) are excluded because the equilibrium effect in the present query is not defined for them.